% !PRINTED 540
\weberpage{540}
%<<CONT
genügen, und ist $\psi(x_1, x_2, \ldots)$ eine ganze Function mit ganz\-
zahligen rationalen Co\"efficienten, so ist auch
\weq{3}{\psi(\xi_1, \xi_2, \ldots) \equiv \psi(\eta_1, \eta_2, \ldots) \ (\mathrm{mod}\ \mu).}

\begin{enumerate}
\item[2.] \gesperrt{Ist $m$ die kleinste positive ganze rationale Zahl,
die durch $\mu$ theilbar ist, so sind zwei nach dem
Modul $\mu$ congruente ganze rationale Zahlen auch
nach dem Modul $m$ congruent} (\S.\,138, 13.).
\end{enumerate}

Ein Hauptsatz über die Congruenzen, zu dessen Beweis wir
jetzt schreiten, ist der:

\begin{enumerate}
\item[3.] \gesperrt{Die Anzahl der nach einem Modul $\mu$ incon\-
gruenten ganzen Zahlen eines Körpers $\Omega$ ist
endlich, und zwar gleich der absoluten Norm
von $\mu$.}
\end{enumerate}

Um ihn zu beweisen, nehmen wir eine Basis des Func\-
tionals $\mu$ an, und zwar wählen wir gerade die im \S.\,146, (3)
bestimmte specielle Basis:
\weq{4}{
\begin{array}{l}
\alpha_1 = a_{1,1}\,\omega_1 \\
\alpha_2 = a_{1,2}\,\omega_1 + a_{2,2}\,\omega_2 \\
\cdots\cdots\cdots\cdots\cdots\cdots\cdots\cdots\cdots\cdots \\
\alpha_n = a_{1,n}\,\omega_1 + a_{2,n}\,\omega_2 + \cdots + a_{n,n}\,\omega_n,
\end{array}
}
worin $\omega_1, \omega_2, \ldots, \omega_n$ eine Basis von $\mathfrak{o}$ ist.

Jede ganze Zahl $\eta$ in $\Omega$ lässt sich dann in der Form dar\-
stellen:
\weq{5}{\eta = y_1\omega_1 + y_2\omega_2 + \cdots + y_n\omega_n,}
wenn die $y_1, y_2, \ldots, y_n$ ganze rationale Zahlen sind. Damit ver\-
binden wir das Zahlensystem
\weq{6}{\xi = x_1\omega_1 + x_2\omega_2 + \cdots + x_n\omega_n,}
in dem wir $x_1$ ein volles Restsystem nach dem Modul $a_{1,1}$,
$x_2$ nach dem Modul $a_{2,2} \ldots$, $x_n$ nach dem Modul $a_{n,n}$ durch\-
laufen lassen. Die Anzahl der so gebildeten Zahlen $\xi$ beträgt
$a_{1,1}a_{2,2}\ldots a_{n,n}$, eine Zahl, die nach \S.\,147, (12) gleich der
absoluten Norm von $\mu$ ist.

Wir wollen nachweisen, dass jede Zahl $\eta$ einer und nur
einer dieser Zahlen $\xi$ nach dem Modul $\mu$ congruent ist.

Wenn $\eta - \xi$ durch $\mu$ theilbar sein soll, so muss nach der
Bedeutung der Basis diese Differenz in der Form
\weq{7}{\eta - \xi = h_1\alpha_1 + h_2\alpha_2 + \cdots + h_n\alpha_n}
%CONT>>
